% Options for packages loaded elsewhere
% Options for packages loaded elsewhere
\PassOptionsToPackage{unicode}{hyperref}
\PassOptionsToPackage{hyphens}{url}
%
\documentclass[
  12pt,
]{scrbook}
\usepackage{xcolor}
\usepackage{amsmath,amssymb}
\setcounter{secnumdepth}{5}
\usepackage{iftex}
\ifPDFTeX
  \usepackage[T1]{fontenc}
  \usepackage[utf8]{inputenc}
  \usepackage{textcomp} % provide euro and other symbols
\else % if luatex or xetex
  \usepackage{unicode-math} % this also loads fontspec
  \defaultfontfeatures{Scale=MatchLowercase}
  \defaultfontfeatures[\rmfamily]{Ligatures=TeX,Scale=1}
\fi
\usepackage{lmodern}
\ifPDFTeX\else
  % xetex/luatex font selection
\fi
% Use upquote if available, for straight quotes in verbatim environments
\IfFileExists{upquote.sty}{\usepackage{upquote}}{}
\IfFileExists{microtype.sty}{% use microtype if available
  \usepackage[]{microtype}
  \UseMicrotypeSet[protrusion]{basicmath} % disable protrusion for tt fonts
}{}
\makeatletter
\@ifundefined{KOMAClassName}{% if non-KOMA class
  \IfFileExists{parskip.sty}{%
    \usepackage{parskip}
  }{% else
    \setlength{\parindent}{0pt}
    \setlength{\parskip}{6pt plus 2pt minus 1pt}}
}{% if KOMA class
  \KOMAoptions{parskip=half}}
\makeatother
% Make \paragraph and \subparagraph free-standing
\makeatletter
\ifx\paragraph\undefined\else
  \let\oldparagraph\paragraph
  \renewcommand{\paragraph}{
    \@ifstar
      \xxxParagraphStar
      \xxxParagraphNoStar
  }
  \newcommand{\xxxParagraphStar}[1]{\oldparagraph*{#1}\mbox{}}
  \newcommand{\xxxParagraphNoStar}[1]{\oldparagraph{#1}\mbox{}}
\fi
\ifx\subparagraph\undefined\else
  \let\oldsubparagraph\subparagraph
  \renewcommand{\subparagraph}{
    \@ifstar
      \xxxSubParagraphStar
      \xxxSubParagraphNoStar
  }
  \newcommand{\xxxSubParagraphStar}[1]{\oldsubparagraph*{#1}\mbox{}}
  \newcommand{\xxxSubParagraphNoStar}[1]{\oldsubparagraph{#1}\mbox{}}
\fi
\makeatother

\usepackage{color}
\usepackage{fancyvrb}
\newcommand{\VerbBar}{|}
\newcommand{\VERB}{\Verb[commandchars=\\\{\}]}
\DefineVerbatimEnvironment{Highlighting}{Verbatim}{commandchars=\\\{\}}
% Add ',fontsize=\small' for more characters per line
\usepackage{framed}
\definecolor{shadecolor}{RGB}{241,243,245}
\newenvironment{Shaded}{\begin{snugshade}}{\end{snugshade}}
\newcommand{\AlertTok}[1]{\textcolor[rgb]{0.68,0.00,0.00}{#1}}
\newcommand{\AnnotationTok}[1]{\textcolor[rgb]{0.37,0.37,0.37}{#1}}
\newcommand{\AttributeTok}[1]{\textcolor[rgb]{0.40,0.45,0.13}{#1}}
\newcommand{\BaseNTok}[1]{\textcolor[rgb]{0.68,0.00,0.00}{#1}}
\newcommand{\BuiltInTok}[1]{\textcolor[rgb]{0.00,0.23,0.31}{#1}}
\newcommand{\CharTok}[1]{\textcolor[rgb]{0.13,0.47,0.30}{#1}}
\newcommand{\CommentTok}[1]{\textcolor[rgb]{0.37,0.37,0.37}{#1}}
\newcommand{\CommentVarTok}[1]{\textcolor[rgb]{0.37,0.37,0.37}{\textit{#1}}}
\newcommand{\ConstantTok}[1]{\textcolor[rgb]{0.56,0.35,0.01}{#1}}
\newcommand{\ControlFlowTok}[1]{\textcolor[rgb]{0.00,0.23,0.31}{\textbf{#1}}}
\newcommand{\DataTypeTok}[1]{\textcolor[rgb]{0.68,0.00,0.00}{#1}}
\newcommand{\DecValTok}[1]{\textcolor[rgb]{0.68,0.00,0.00}{#1}}
\newcommand{\DocumentationTok}[1]{\textcolor[rgb]{0.37,0.37,0.37}{\textit{#1}}}
\newcommand{\ErrorTok}[1]{\textcolor[rgb]{0.68,0.00,0.00}{#1}}
\newcommand{\ExtensionTok}[1]{\textcolor[rgb]{0.00,0.23,0.31}{#1}}
\newcommand{\FloatTok}[1]{\textcolor[rgb]{0.68,0.00,0.00}{#1}}
\newcommand{\FunctionTok}[1]{\textcolor[rgb]{0.28,0.35,0.67}{#1}}
\newcommand{\ImportTok}[1]{\textcolor[rgb]{0.00,0.46,0.62}{#1}}
\newcommand{\InformationTok}[1]{\textcolor[rgb]{0.37,0.37,0.37}{#1}}
\newcommand{\KeywordTok}[1]{\textcolor[rgb]{0.00,0.23,0.31}{\textbf{#1}}}
\newcommand{\NormalTok}[1]{\textcolor[rgb]{0.00,0.23,0.31}{#1}}
\newcommand{\OperatorTok}[1]{\textcolor[rgb]{0.37,0.37,0.37}{#1}}
\newcommand{\OtherTok}[1]{\textcolor[rgb]{0.00,0.23,0.31}{#1}}
\newcommand{\PreprocessorTok}[1]{\textcolor[rgb]{0.68,0.00,0.00}{#1}}
\newcommand{\RegionMarkerTok}[1]{\textcolor[rgb]{0.00,0.23,0.31}{#1}}
\newcommand{\SpecialCharTok}[1]{\textcolor[rgb]{0.37,0.37,0.37}{#1}}
\newcommand{\SpecialStringTok}[1]{\textcolor[rgb]{0.13,0.47,0.30}{#1}}
\newcommand{\StringTok}[1]{\textcolor[rgb]{0.13,0.47,0.30}{#1}}
\newcommand{\VariableTok}[1]{\textcolor[rgb]{0.07,0.07,0.07}{#1}}
\newcommand{\VerbatimStringTok}[1]{\textcolor[rgb]{0.13,0.47,0.30}{#1}}
\newcommand{\WarningTok}[1]{\textcolor[rgb]{0.37,0.37,0.37}{\textit{#1}}}

\usepackage{longtable,booktabs,array}
\usepackage{calc} % for calculating minipage widths
% Correct order of tables after \paragraph or \subparagraph
\usepackage{etoolbox}
\makeatletter
\patchcmd\longtable{\par}{\if@noskipsec\mbox{}\fi\par}{}{}
\makeatother
% Allow footnotes in longtable head/foot
\IfFileExists{footnotehyper.sty}{\usepackage{footnotehyper}}{\usepackage{footnote}}
\makesavenoteenv{longtable}
\usepackage{graphicx}
\makeatletter
\newsavebox\pandoc@box
\newcommand*\pandocbounded[1]{% scales image to fit in text height/width
  \sbox\pandoc@box{#1}%
  \Gscale@div\@tempa{\textheight}{\dimexpr\ht\pandoc@box+\dp\pandoc@box\relax}%
  \Gscale@div\@tempb{\linewidth}{\wd\pandoc@box}%
  \ifdim\@tempb\p@<\@tempa\p@\let\@tempa\@tempb\fi% select the smaller of both
  \ifdim\@tempa\p@<\p@\scalebox{\@tempa}{\usebox\pandoc@box}%
  \else\usebox{\pandoc@box}%
  \fi%
}
% Set default figure placement to htbp
\def\fps@figure{htbp}
\makeatother





\setlength{\emergencystretch}{3em} % prevent overfull lines

\providecommand{\tightlist}{%
  \setlength{\itemsep}{0pt}\setlength{\parskip}{0pt}}



 


\usepackage{xcolor}
% define colors
\newcommand{\red}[1]{\color[RGB]{185,64,71}{#1}}
\newcommand{\green}[1]{\color[RGB]{0, 139, 69}{#1}}
\definecolor{red}{RGB}{185,64,71}

% Configure hyperref for expanded bookmarks
% Note that this only works for Adobe Acrobat Reader
\usepackage{hyperref}
\hypersetup{
  bookmarksopen = true, % show TOC with all the subtrees expanded
  bookmarksopenlevel = 2 % level to which bookmarks are open
}

%% bullet list
\usepackage{enumitem}
\setlist[itemize, 1]{
    itemsep = 0.5\baselineskip,
    parsep = 0.2\baselineskip}
\setlist[itemize, 2]{
    itemsep = 0.2\baselineskip
}
\renewcommand{\tightlist}{\setlength{\itemsep}{0.5\baselineskip}\setlength{\parskip}{0.2\baselineskip}}

%% set font
\usepackage{fontspec}
\setmainfont{XCharter}
\setmonofont{Consolas}
\usepackage{unicode-math}
\setmathfont{XCharter-Math.otf} % do NOT delete the .otf extension

%% define custom unicode characters
\usepackage{newunicodechar}
%% for symbols that are not included in the main font, we can use a more symbol robust font to render them, here we use Apple Symbols
\newfontfamily\symbolfont{DejaVu Sans} % Sans is more readable for symbols, rounder and thicker than Serif; Sans Mono is a bit more compact
\newunicodechar{→}{{\symbolfont →}}
\newunicodechar{≈}{{\symbolfont ≈}}
\newunicodechar{⨯}{{\symbolfont ⨯}}
\newunicodechar{♂}{{\symbolfont ♂}}
\newunicodechar{♀}{{\symbolfont ♀}}

%% support for Chinese
\usepackage{xeCJK}
\setCJKmainfont{LXGW WenKai}
\setCJKsansfont{Microsoft YaHei}
\setCJKmonofont{Maple Mono NF CN}

\usepackage{changepage}   % for the adjustwidth environment
\usepackage[all]{nowidow} % prevent widow/orphan lines
\pretolerance = 1000               % Relaxes parameters for line breaks
\tolerance    = 2000               % Relaxes parameters for line breaks


%% table typesetting
\usepackage{threeparttable, adjustbox}
\usepackage{booktabs}
\renewcommand{\arraystretch}{1.3} % increase row spacing
\usepackage{float} % for [H] placement specifier
% Make all tables use [H] placement by default
% Still takes placement, eg [htbp]; defaults to [H] if no placement is specified
\makeatletter
\renewenvironment{table}[1][H]{%
  \@float{table}[#1]%
}{%
  \end@float
}
\makeatother
\usepackage[labelfont={bf}, labelsep=space, skip=4pt]{caption}
\usepackage{siunitx}
% define new column types
\newcommand\mc[1]{\multicolumn{1}{c}{#1}} % center table header
\usepackage{dcolumn}
\newcolumntype{d}[1]{D{.}{.}{#1}} % align at decimal point

%% Math symbols
\usepackage{amsmath, amssymb, textcomp}

%% define TeX macros

% bold face capital letter 大写加粗
\newcommand{\bA}{\boldsymbol{A}}
\newcommand{\bB}{\boldsymbol{B}}
\newcommand{\bC}{\boldsymbol{C}}
\newcommand{\bD}{\boldsymbol{D}}
\newcommand{\bE}{\boldsymbol{E}}
\newcommand{\bF}{\boldsymbol{F}}
\newcommand{\bG}{\boldsymbol{G}}
\newcommand{\bH}{\boldsymbol{H}}
\newcommand{\bI}{\boldsymbol{I}}
\newcommand{\bJ}{\boldsymbol{J}}
\newcommand{\bK}{\boldsymbol{K}}
\newcommand{\bL}{\boldsymbol{L}}
\newcommand{\bM}{\boldsymbol{M}}
\newcommand{\bN}{\boldsymbol{N}}
\newcommand{\bP}{\boldsymbol{P}}
\newcommand{\bQ}{\boldsymbol{Q}}
\newcommand{\bR}{\boldsymbol{R}}
\newcommand{\bS}{\boldsymbol{S}}
\newcommand{\bT}{\boldsymbol{T}}
\newcommand{\bU}{\boldsymbol{U}}
\newcommand{\bV}{\boldsymbol{V}}
\newcommand{\bW}{\boldsymbol{W}}
\newcommand{\bX}{\boldsymbol{X}}
\newcommand{\bY}{\boldsymbol{Y}}
\newcommand{\bZ}{\boldsymbol{Z}}

% bold face small letter 小写加粗
\newcommand{\ba}{\boldsymbol{a}}
\newcommand{\bb}{\boldsymbol{b}}
\newcommand{\bc}{\boldsymbol{c}}
\newcommand{\bd}{\boldsymbol{d}}
\newcommand{\be}{\boldsymbol{e}}
\newcommand{\bg}{\boldsymbol{g}}
\newcommand{\bh}{\boldsymbol{h}}
\newcommand{\bi}{\boldsymbol{i}}
\newcommand{\bm}{\boldsymbol{m}}
\newcommand{\bp}{\boldsymbol{p}}
\newcommand{\bq}{\boldsymbol{q}}
\newcommand{\br}{\boldsymbol{r}}
\newcommand{\bs}{\boldsymbol{s}}
\newcommand{\bt}{\boldsymbol{t}}
\newcommand{\bu}{\boldsymbol{u}}
\newcommand{\bv}{\boldsymbol{v}}
\newcommand{\bw}{\boldsymbol{w}}
\newcommand{\bx}{\boldsymbol{x}}
\newcommand{\by}{\boldsymbol{y}}

% math calligraphic font 花体
\newcommand{\Acal}{\mathcal{A}}
\newcommand{\Bcal}{\mathcal{B}}
\newcommand{\Ccal}{\mathcal{C}}
\newcommand{\Dcal}{\mathcal{D}}
\newcommand{\Ecal}{\mathcal{E}}
\newcommand{\Fcal}{\mathcal{F}}
\newcommand{\Gcal}{\mathcal{G}}
\newcommand{\Hcal}{\mathcal{H}}
\newcommand{\Ical}{\mathcal{I}}
\newcommand{\Jcal}{\mathcal{J}}
\newcommand{\Kcal}{\mathcal{K}}
\newcommand{\Lcal}{\mathcal{L}}
\newcommand{\Mcal}{\mathcal{M}}
\newcommand{\Ncal}{\mathcal{N}}
\newcommand{\Ocal}{\mathcal{O}}
\newcommand{\Pcal}{\mathcal{P}}
\newcommand{\Qcal}{\mathcal{Q}}
\newcommand{\Rcal}{\mathcal{R}}
\newcommand{\Scal}{\mathcal{S}}
\newcommand{\Tcal}{\mathcal{T}}
\newcommand{\Ucal}{\mathcal{U}}
\newcommand{\Vcal}{\mathcal{V}}
\newcommand{\Wcal}{\mathcal{W}}
\newcommand{\Xcal}{\mathcal{X}}
\newcommand{\Ycal}{\mathcal{Y}}
\newcommand{\Zcal}{\mathcal{Z}}

% blackboard face 黑板粗体
\newcommand{\A}{\mathbb{A}}
\newcommand{\B}{\mathbb{B}}
\newcommand{\C}{\mathbb{C}}
\newcommand{\D}{\mathbb{D}}
\newcommand{\E}{\mathbb{E}}
\newcommand{\F}{\mathbb{F}}
\newcommand{\G}{\mathbb{G}}
\renewcommand{\H}{\mathbb{H}}
\newcommand{\I}{\mathbb{I}}
\newcommand{\J}{\mathbb{J}}
\newcommand{\K}{\mathbb{K}}
\renewcommand{\L}{\mathbb{L}}
\newcommand{\M}{\mathbb{M}}
\newcommand{\N}{\mathbb{N}}
\renewcommand{\O}{\mathbb{O}}
\renewcommand{\P}{\mathbb{P}}
\newcommand{\Q}{\mathbb{Q}}
\newcommand{\R}{\mathbb{R}}
\renewcommand{\S}{\mathbb{S}}
\newcommand{\T}{\mathbb{T}}
\newcommand{\U}{\mathbb{U}}
\newcommand{\V}{\mathbb{V}}
\newcommand{\W}{\mathbb{W}}
\newcommand{\X}{\mathbb{X}}
\newcommand{\Y}{\mathbb{Y}}
\newcommand{\Z}{\mathbb{Z}}

% bold face uppercase Greeks
\newcommand{\bAlpha}{\boldsymbol{\Alpha}}
\newcommand{\bBeta}{\boldsymbol{\Beta}}
\newcommand{\bDelta}{\boldsymbol{\Delta}}
\newcommand{\bEta}{\boldsymbol{\eta}}
\newcommand{\bGamma}{\boldsymbol{\Gamma}}
\newcommand{\bLambda}{\boldsymbol{\Lambda}}
\newcommand{\bOmega}{\boldsymbol{\Omega}}
\newcommand{\bPhi}{\boldsymbol{\Phi}}
\newcommand{\bPi}{\boldsymbol{\Pi}}
\newcommand{\bPsi}{\boldsymbol{\Psi}}
\newcommand{\bSigma}{\boldsymbol{\Sigma}}
\newcommand{\bTau}{\boldsymbol{\tau}}
\newcommand{\bXi}{\boldsymbol{\Xi}}

% bold face lowercase Greeks
\newcommand{\bbeta}{\boldsymbol{\beta}}
\newcommand{\bdelta}{\boldsymbol{\delta}}
\newcommand{\blambda}{\boldsymbol{\lambda}}
\newcommand{\bmu}{\boldsymbol{\mu}}
\newcommand{\bphi}{\boldsymbol{\phi}}
\newcommand{\bpi}{\boldsymbol{\pi}}
\newcommand{\bpsi}{\boldsymbol{\psi}}
\newcommand{\brho}{\boldsymbol{\rho}}
\newcommand{\bsigma}{\boldsymbol{\sigma}}
\newcommand{\btau}{\boldsymbol{\tau}}
\newcommand{\bxi}{\boldsymbol{\xi}}

% Functions (single argument)
\newcommand{\curl}[1]{\left\lbrace #1 \right\rbrace}
\newcommand{\floor}[1]{\left\lfloor #1 \right\rfloor}
\newcommand{\ceil}[1]{\left\lceil #1 \right\rceil}
\newcommand{\abs}[1]{\left\lvert #1 \right\rvert}
\newcommand{\norm}[1]{\left\lVert #1 \right\rVert}

% Other operators
\newcommand{\var}{\mathrm{Var}}
\newcommand{\cov}{\mathrm{Cov}}
\newcommand{\cor}{\mathrm{Corr}}
\newcommand{\diag}{\mathrm{diag}}
\newcommand{\rank}{\mathrm{rank}}
\newcommand{\vecc}{\mathrm{vec}}
\newcommand{\indep}{\perp \! \! \! \perp}

% Remove "0." prefix from subsection and subsubsection numbering
\makeatletter
\renewcommand{\thesubsection}{%
  \ifnum\value{section}=0
    \arabic{subsection}%
  \else
    \thesection.\arabic{subsection}%
  \fi
}
\renewcommand{\thesubsubsection}{%
  \ifnum\value{subsection}=0
    \ifnum\value{section}=0
      \arabic{subsubsection}%
    \else
      \thesection.\arabic{subsubsection}%
    \fi
  \else
    \ifnum\value{section}=0
      \arabic{subsection}.\arabic{subsubsection}%
    \else
      \thesection.\arabic{subsection}.\arabic{subsubsection}%
    \fi
  \fi
}
\makeatother
\usepackage{xcolor}
\usepackage{colortbl}
\makeatletter
\@ifpackageloaded{tcolorbox}{}{\usepackage[skins,breakable]{tcolorbox}}
\@ifpackageloaded{fontawesome5}{}{\usepackage{fontawesome5}}
\definecolor{quarto-callout-color}{HTML}{909090}
\definecolor{quarto-callout-note-color}{HTML}{0758E5}
\definecolor{quarto-callout-important-color}{HTML}{CC1914}
\definecolor{quarto-callout-warning-color}{HTML}{EB9113}
\definecolor{quarto-callout-tip-color}{HTML}{00A047}
\definecolor{quarto-callout-caution-color}{HTML}{FC5300}
\definecolor{quarto-callout-color-frame}{HTML}{acacac}
\definecolor{quarto-callout-note-color-frame}{HTML}{4582ec}
\definecolor{quarto-callout-important-color-frame}{HTML}{d9534f}
\definecolor{quarto-callout-warning-color-frame}{HTML}{f0ad4e}
\definecolor{quarto-callout-tip-color-frame}{HTML}{02b875}
\definecolor{quarto-callout-caution-color-frame}{HTML}{fd7e14}
\makeatother
\makeatletter
\@ifpackageloaded{caption}{}{\usepackage{caption}}
\AtBeginDocument{%
\ifdefined\contentsname
  \renewcommand*\contentsname{Table of contents}
\else
  \newcommand\contentsname{Table of contents}
\fi
\ifdefined\listfigurename
  \renewcommand*\listfigurename{List of Figures}
\else
  \newcommand\listfigurename{List of Figures}
\fi
\ifdefined\listtablename
  \renewcommand*\listtablename{List of Tables}
\else
  \newcommand\listtablename{List of Tables}
\fi
\ifdefined\figurename
  \renewcommand*\figurename{Figure}
\else
  \newcommand\figurename{Figure}
\fi
\ifdefined\tablename
  \renewcommand*\tablename{Table}
\else
  \newcommand\tablename{Table}
\fi
}
\@ifpackageloaded{float}{}{\usepackage{float}}
\floatstyle{ruled}
\@ifundefined{c@chapter}{\newfloat{codelisting}{h}{lop}}{\newfloat{codelisting}{h}{lop}[chapter]}
\floatname{codelisting}{Listing}
\newcommand*\listoflistings{\listof{codelisting}{List of Listings}}
\usepackage{amsthm}
\theoremstyle{plain}
\newtheorem{theorem}{Theorem}[section]
\theoremstyle{definition}
\newtheorem{example}{Example}[section]
\theoremstyle{remark}
\AtBeginDocument{\renewcommand*{\proofname}{Proof}}
\newtheorem*{remark}{Remark}
\newtheorem*{solution}{Solution}
\newtheorem{refremark}{Remark}[section]
\newtheorem{refsolution}{Solution}[section]
\makeatother
\makeatletter
\makeatother
\makeatletter
\@ifpackageloaded{caption}{}{\usepackage{caption}}
\@ifpackageloaded{subcaption}{}{\usepackage{subcaption}}
\makeatother
\usepackage{bookmark}
\IfFileExists{xurl.sty}{\usepackage{xurl}}{} % add URL line breaks if available
\urlstyle{same}
\hypersetup{
  pdftitle={Theoretical Framework of Linear Regression},
  hidelinks,
  pdfcreator={LaTeX via pandoc}}


\title{Theoretical Framework of Linear Regression}
\author{}
\date{}
\begin{document}
\frontmatter
\maketitle


\mainmatter
\begin{tcolorbox}[enhanced jigsaw, colback=white, colframe=quarto-callout-note-color-frame, leftrule=.75mm, bottomrule=.15mm, opacityback=0, rightrule=.15mm, arc=.35mm, toprule=.15mm, breakable, left=2mm]

🎯 \textbf{Study Objectives}

\begin{itemize}
\tightlist
\item
  Understand the key assumptions underlying the linear regression model.
\item
  Understand the theoretical framework of linear regression, including
  the OLS estimator, its properties, and its distribution.
\item
  Central Limit Theorem (CLT)

  \begin{itemize}
  \tightlist
  \item
    Apply the CLT to approximate the distribution of sample mean and
    calculating probabilities using distribution table.
  \item
    Understand the implications of CLT for the sampling distribution of
    the OLS estimator.
  \end{itemize}
\item
  Distinguish the difference between exact distribution and sampling
  distribution of the OLS estimator.
\item
  Understand R squared and adjusted R squared as measures of fit.
\item
  Be able to identify violations of the linear regression assumptions
  using diagnostic plots.
\item
  Understand how to address heteroskedasticity.
\end{itemize}

\end{tcolorbox}

In this section, we will discuss the theoretical framework of linear
regression, including the assumptions underlying the linear regression
model and essential hypothesis tests.

Q: Why the assumptions matter?

A: In short, violation of the assumptions may lead to biased or
inefficient estimates, incorrect inferences, and unreliable predictions.

\begin{center}\rule{0.5\linewidth}{0.5pt}\end{center}

\section{Assumptions of Linear
Regression}\label{assumptions-of-linear-regression}

Q: What are the assumptions of linear regression?

A: The \textbf{Gauss-Markov Assumptions} ensure that the Ordinary Least
Squares (OLS) estimator is the Best Linear Unbiased Estimator (BLUE).

\begin{enumerate}
\def\labelenumi{\arabic{enumi}.}
\item
  \textbf{Linearity}: The relationship between the independent
  variable(s) and the dependent variable is linear in parameters.
\item
  \textbf{No perfect multicollinearity}: The independent variables are
  NOT {perfectly correlated.}
\item
  \textbf{Error has zero conditional mean}
\item
  \textbf{Error is homoskedastic}

  \begin{itemize}
  \tightlist
  \item
    Homoskedasticity: The variance of the error term is constant across
    all levels of the independent variable(s).
  \end{itemize}

  {In time series settings (we won't cover in this course), we also
  assume:}

  \begin{itemize}
  \tightlist
  \item
    {No autocorrelation: The error terms are not correlated with each
    other.}
  \end{itemize}
\item
  \textbf{Residuals are normally distributed}

  This assumption is useful for conducting hypothesis tests and
  constructing confidence intervals.
\item
  \textbf{Large outliers are unlikely}

  This assumption requires that observations with values of \(X_i\),
  \(Y_i\), or both that are far outside the usual range of the data are
  unlikely.

  Large outliers can make OLS regression results misleading as OLS are
  sensitive to outliers.

  Mathematically, this assumption can be stated as: \(X\) and \(Y\) have
  \textbf{finite kurtosis}.
\end{enumerate}

\section{Regression Framework and OLS
Estimator}\label{regression-framework-and-ols-estimator}

Multiple linear regression models is defined as

\[
Y_i = \beta_0 + \beta_1 X_{i1} + \beta_2 X_{i2} + ... + \beta_k X_{ik} + \varepsilon_i
\]

where

\begin{itemize}
\tightlist
\item
  \(i=1,\ldots, n\) indexes the observations,
\item
  \(Y_i\) is the dependent variable, and
\item
  \(X_{ij}\) represents the \(j\)-th independent variable for
  observation \(i\), with \(j = 1, \ldots, k\) and \(k\) denoting the
  total number of independent variables.
\item
  The coefficients \(\beta_j\) are to be estimated, and
\item
  \(\varepsilon_i\) is the error term.
\end{itemize}

\textbf{The OLS estimator \(\hat{\beta}^{OLS}\) minimizes the sum of
squared residuals (SSR)}:

\[
\hat{\beta}^{OLS} = \arg\min_{\beta} \sum_{i=1}^{n} e_i^2 
\]

where \(e_i = Y_i - \hat{Y}_i\) is the residual for observation \(i\).

\(\hat{Y}_i\) is the predicted value of \(Y_i\) given the estimated
coefficients.

\[
\hat{\beta}^{OLS} = \arg\min_{\beta} \sum_{i=1}^{n} (Y_i - \beta_0 - \sum_{j=1}^{k} \beta_j X_{ij})^2
\]

\(\hat{\beta}^{OLS}\) can be obtained by the first order condition
(setting the derivative of SSR with respect to each \(\beta_j\) to zero
and solving the resulting system of equations).

Read
\href{https://github.com/my1396/FIN5005-Fall2024/blob/a4ec0c87c6bcbd463b4da1ad78878b1fe329cda4/images/matrix_OLS_NYU_notes.pdf}{OLS
in Matrix Form} for more details. We won't cover the OLS derivation in
this course.

Here we give the direct formula for the OLS estimator:

\[
\hat\beta = \left(\sum_{i=1}^n X_iX_i'\right)^{-1} \left(\sum_{i=1}^n X_iY_i\right)
\]

where

\[
X_i = \begin{pmatrix}1 \\ X_{i1} \\ X_{i2} \\ \vdots \\ X_{ik}\end{pmatrix}
\]

is a \((k+1) \times 1\) vector of independent variables including the
intercept for observation \(i\).

\textbf{Properties of OLS estimator:}

\begin{itemize}
\tightlist
\item
  \textbf{Unbiasedness}: \(E[\hat{\beta}^{OLS}] = \beta\) (if
  Gauss-Markov assumptions hold)
\item
  \textbf{Efficiency}: Among all linear unbiased estimators, OLS has the
  smallest variance
\item
  \textbf{Consistency}: As sample size \(n \to \infty\),
  \(\hat{\beta}^{OLS} \to \beta\) in probability
\end{itemize}

\subsection{Exact Distribution of OLS
Estimator}\label{exact-distribution-of-ols-estimator}

Under the assumption of normally distributed errors,
\(\boldsymbol{\varepsilon} \vert \boldsymbol{X} \sim N(\boldsymbol{0}, \sigma^2\boldsymbol{I})\),
we have in matrix form

\begin{equation}\phantomsection\label{eq-beta-distribution}{
\hat{\boldsymbol{\beta}} \sim N(\boldsymbol{\beta}, \sigma^2 (\boldsymbol{X}'\boldsymbol{X})^{-1}).
}\end{equation}

where

\begin{itemize}
\tightlist
\item
  \(\sigma^2 = \var (\varepsilon_i \mid \bX)\) is the variance of the
  error term.
\item
  \(\bX\) is the matrix of independent variables. \[
  \bX = \begin{pmatrix}
     X_1' \\ X_2' \\ \vdots \\ X_n'
  \end{pmatrix}_{n \times (k+1)}
  \]
\end{itemize}

Equation~\ref{eq-beta-distribution} is called the \textbf{exact
(finite-sample) distribution} of the OLS estimator.

\(\var(\hat{\bbeta}) = \sigma^2 (\boldsymbol{X}'\boldsymbol{X})^{-1}\)
is called the \textbf{variance-covariance matrix} of the OLS estimator,
often denoted as \(V_{\hat\beta}.\)

We estimate \(\sigma^2\) with \(\hat{\sigma}^2:\) \[
\hat{\sigma}^2 = \frac{\sum_{i=1}^ne_i^2}{n-k-1}.
\]

The standard errors of \(\hat{\beta}_j\), \(j=1,\ldots,k,\) are given by
the square root of the \((j,j)\) element of the variance-covariance
matrix, \(V_{\hat\beta}\).

The individual coefficient estimates \(\hat{\beta}_j\) are normally
distributed:

\[
\hat{\beta}_j \sim N(\beta_j, \sigma^2 v_{jj})
\]

\section{Sampling Distribution of OLS
Estimator}\label{sampling-distribution-of-ols-estimator}

Becausethe OLS estimator \(\hat\beta_0, \hat\beta_1, \ldots\) are
computed from a randomly drawn sample of data, the estimators themselves
are random variables with a probability distribution.

The sampling distribution describes the values they could take over
different possible random samples. In small samples, theses sampling
distributions are complicated, but in large samples, they are normal
because of the Central Limit Theorem (CLT).

\begin{tcolorbox}[enhanced jigsaw, colback=white, coltitle=black, title=\textcolor{quarto-callout-tip-color}{\faLightbulb}\hspace{0.5em}{Sampling Distribution vs.~Exact Distribution}, bottomtitle=1mm, opacityback=0, rightrule=.15mm, left=2mm, colbacktitle=quarto-callout-tip-color!10!white, colframe=quarto-callout-tip-color-frame, leftrule=.75mm, bottomrule=.15mm, opacitybacktitle=0.6, toprule=.15mm, titlerule=0mm, arc=.35mm, breakable, toptitle=1mm]

\begin{itemize}
\tightlist
\item
  The distribution depends on the sample data \((X_i, Y_i).\)
\item
  Based on Central Limit Theorem (CLT), we can establish the sampling
  distribution of \(\hat{\bbeta}\) when sample size is large, i.e.,
  \(n\to\infty\).
\end{itemize}

\end{tcolorbox}

\subsection{Central Limit Theorem
(CLT)}\label{central-limit-theorem-clt}

\begin{theorem}[]\protect\hypertarget{thm-clt}{}\label{thm-clt}

Suppose that \(X_1, \ldots, X_n\) is an independent and identically
distributed (iid) sequence with a finite mean \(\mathbb{E}(X_i)=\mu\)
and variance \(\text{Var}(X_i)=\sigma^2\), where \(0<\sigma^2<\infty\).

Define a sample mean: \(\overline{X}=\frac{1}{n}\sum_{i=1}^n X_i.\)

The {\textbf{Lindeberg-Lévy CLT}} states that

\[
\frac{\overline{X}-\mathbb{E}[\overline{X}]}{\sqrt{\text{Var}[\overline{X}]}}
= \frac{\overline{X}-\mu}{\sqrt{\sigma^2/n}}
= {\color[HTML]{00CC66} \sqrt{n} \cdot \frac{\overline{X}-\mu}{\sigma} }
\xrightarrow{d} N(0,1)
\] where \(\xrightarrow{d}\) denotes convergence in distribution.

Alternative expressions of CLT: \[
\begin{split}
\sqrt{n} (\overline{X}-\mu)  & \xrightarrow{d} N(0,\sigma^2) \\
\red{\overline{X}} & \red{\xrightarrow{d} N(\mu, \frac{\sigma^2}{n})}
\end{split}
\]

\end{theorem}

Q: Why does CLT matter?

A:

\begin{itemize}
\item
  The central limit theorem implies that if \(N\) is large
  \textbf{enough}, we can \textbf{approximate} the distribution of
  \(\overline{X}\) by the {\textbf{standard normal distribution}} with
  mean \(\mu\) and variance \(\sigma^2 / N\) \textbf{regardless of the
  underlying distribution of \(X.\)}
\item
  The scaling by {\(n\)} is crucial. \[
  \sigma_{\overline{X}}^2 = \text{Var}[\overline{X}] = \frac{\sigma^2}{n}
  \] This means as sample size increases, the sample variance shrinks at
  the rate of \(1/n\).

  \begin{tcolorbox}[enhanced jigsaw, colback=white, coltitle=black, title=\textcolor{quarto-callout-tip-color}{\faLightbulb}\hspace{0.5em}{Tip}, bottomtitle=1mm, opacityback=0, rightrule=.15mm, left=2mm, colbacktitle=quarto-callout-tip-color!10!white, colframe=quarto-callout-tip-color-frame, leftrule=.75mm, bottomrule=.15mm, opacitybacktitle=0.6, toprule=.15mm, titlerule=0mm, arc=.35mm, breakable, toptitle=1mm]

  \textbf{Interpretation:} As the sample size \(n\) increases,the mean
  becomes more concentrated around the true population mean \(\mu\).

  \end{tcolorbox}
\item
  How large must \(n\) be for the distribution of \(\overline{X}\) to be
  approximately normal?

  There is no definitive answer, but a common rule of thumb is that
  \(n \geq 30\) is often sufficient for the CLT to provide a good
  approximation.

  Because the distribution of \(\overline{X}\) approaches the normal
  distribution as \(n\) increases, \(\overline{X}\) is said to have an
  {\textbf{asymptotic normal distribution}}.
\end{itemize}

\subsection{CLT Example}\label{clt-example}

The amount of time customers spend in a grocery store is a random
variable with mean \(\mu = 40\) minutes and standard deviation
\(\sigma = 15\) minutes. Assume each customer's time in the store is
independent of others.

\[
\mathbb{E}[X] = \mu = 40, \quad \text{and} \quad \text{SD}(X) = \sigma = 15
\]

Consider the following probabilities:

\begin{example}[]\protect\hypertarget{exm-clt1}{}\label{exm-clt1}

Assuming \(X\) is normally distributed, what is the probability that a
randomly selected customer spends more than 42 minutes in the store,
i.e., compute \(P(X > 42)\)?

Use the normal distribution table to look up the probabilities.

\end{example}

Solution

\phantomsection\label{clt1}
\textbf{Individual Customer}

We define the standard normal variable \(Z\) as:

\[
Z = \frac{X - \mu}{\sigma} \sim N(0,1)
\]

where \(\mu = 40\) and \(\sigma = 15\).

To compute \(P(X > 42)\), we convert to the standard normal:

\[
P(X > 42) = P(\frac{X-40}{15} > \frac{42-40}{15}) =  P(Z > 0.1333)
\]

By symmetry of the normal distribution:

\[P(Z > 0.1333) = P(Z < -0.1333) = \Phi(-0.1333) = 0.4483\]

where \(\Phi(\cdot)\) denotes the cumulative distribution function (CDF)
of the standard normal distribution.

\textbf{Interpretation:} There is approximately a \textbf{44.83\%}
probability that a randomly selected customer will spend more than 42
minutes in the store.

Note that the \textbf{specific distribution} for \(X\), i.e., normality
here, is \textbf{required} to compute the probability for a single
customer.

\begin{example}[]\protect\hypertarget{exm-clt2}{}\label{exm-clt2}

Given a random sample of \(n=64\) customers, what is the probability of
the average time spent by the 64 customers exceeds 42 minutes, i.e.,
compute \(P(\overline{X}_{64} > 42)\)?

Hint: apply CLT and justify why CLT can be applied here.

\end{example}

Solution

\phantomsection\label{clt2}
\textbf{Sample Mean for \(n = 64\)}

By the Central Limit Theorem, for \textbf{large \(n\) with finite mean
and variance}, the sampling distribution of the sample mean
\(\overline{X}_{64}\) is approximately normal (regardless the
distribution of \(X\)):

\[
\overline{X}_{64} \sim N (\mu, \frac{\sigma^2}{n})
\]

Plug in \(\mu = 40\) and \(\sigma = 15\):

\[
\sigma_{\overline{X}_{64}} = \frac{\sigma}{\sqrt{n}} = \frac{15}{\sqrt{64}} = 1.875
\]

According to CLT \(\overline{X}_{64}\) is normally distributed:

\[
\overline{X}_{64} \sim N\left(40, 1.875^2\right)
\]

To calculate \(P(\overline{X} > 42)\), we standardize the sample average
to the standard normal variable \(Z\), defined as:

\[
Z = \frac{\overline{X} - \mu}{\sigma/\sqrt{n}} = \frac{\overline{X} - 40}{1.875}
\]

We can rewrite the probability in terms of \(Z\):

\[
P(\overline{X} > 42) = P(\frac{\overline{X} - \mu}{\sigma/\sqrt{n}} > \frac{42 - 40}{1.875}) = P(Z > 1.067)
\]

Using the standard normal distribution table:

\[
P(Z > 1.067) = P(Z < -1.067)
\]

Since the standard normal distribution table is rounded at two decimal
places, we look up \(P(Z < -1.07) = 0.1423\):

\textbf{Interpretation:} There is a \textbf{14.23\%} probability that
the average time spent by a random sample of 64 customers exceeds 42
minutes.

\begin{example}[]\protect\hypertarget{exm-clt3}{}\label{exm-clt3}

Given a random sample of \(n=81\) customers, what is the probability of
the average time spent by the 81 customers exceeds 42 minutes, i.e.,
compute \(P(\overline{X}_{81} > 42)\)?

\end{example}

Solution

\phantomsection\label{clt3}
\textbf{Sample Mean for \(n = 81\)}

\[
\sigma_{\overline{X}_{81}} = \frac{\sigma}{\sqrt{n}} = \frac{15}{\sqrt{81}} = \frac{15}{9} = 1.667
\]

According to CLT \(\overline{X}_{81}\) is normally distributed:

\[
\overline{X}_{81} \sim N(\mu= 40, \sigma_{\overline{X}_{81}} = 1.667^2)
\]

Standardize the sample average

\[
P(\overline{X} > 42) = P(\frac{\overline{X}-\mu}{\sigma_{\overline{X}_{81}}} > \frac{42-40}{1.667})
\]

Rewrite the probability in terms of \(Z\):

\[
z = \frac{42 - 40}{1.667} = 1.20
\]

\[
P(\overline{X} > 42) = P(Z > 1.20) 
\]

Look up the normal distribution table \[
P(Z > 1.20) = P(Z<-1.2) \approx 0.1151
\]

\textbf{Interpretation:} A random sample of 16 customers has a
\textbf{11.51\%} chance of yielding an average time above 42 minutes.

\begin{example}[]\protect\hypertarget{exm-clt4}{}\label{exm-clt4}

Discussion: Compare your results in the three probabilities. Provide a
brief interpretation of how the probabilities differ and why.

\end{example}

Solution

\phantomsection\label{clt4}
\begin{itemize}
\item
  For \textbf{a single customer}, there is a \textbf{44.83\%} chance of
  spending more than 42 minutes.
\item
  For the \textbf{average of 64 customers}, the probability drops to
  \textbf{14.23\%}.
\item
  For the \textbf{average of 81 customers}, it drops further to
  \textbf{11.51\%}.
\item
  \textbf{Why?} As sample size increases, the variability (standard
  deviation) of the sample mean decreases, making it less likely to
  observe extreme values like an average over 42 minutes. This
  demonstrates the Central Limit Theorem and the stabilizing effect of
  larger samples. (See Figure~\ref{fig-clt} for a graphical
  illustration.)

  \begin{itemize}
  \tightlist
  \item
    The individual distribution is wider (larger standard deviation), so
    the chance of extreme values is higher.
  \item
    As sample size increases, the standard error decreases, so the
    distribution of the sample mean becomes narrower.
  \item
    This means it becomes less likely to observe a sample mean far from
    the population mean.
  \end{itemize}
\end{itemize}

\begin{figure}[H]

\centering{

\includegraphics[width=6.6in,height=\textheight,keepaspectratio]{images/CLT.png}

}

\caption{\label{fig-clt}}

\end{figure}%

\begin{tcolorbox}[enhanced jigsaw, colback=white, coltitle=black, title=\textcolor{quarto-callout-tip-color}{\faLightbulb}\hspace{0.5em}{Key Takeaways from CLT}, bottomtitle=1mm, opacityback=0, rightrule=.15mm, left=2mm, colbacktitle=quarto-callout-tip-color!10!white, colframe=quarto-callout-tip-color-frame, leftrule=.75mm, bottomrule=.15mm, opacitybacktitle=0.6, toprule=.15mm, titlerule=0mm, arc=.35mm, breakable, toptitle=1mm]

Generally, as the sample size increases, the \textbf{sampling
distribution of the sample mean} becomes more concentrated around the
population mean.

This reflects a reduction in variability: \textbf{larger samples yield
more precise estimates of the population mean}.

Conversely, smaller sample sizes result in a wider sampling
distribution, indicating greater variability in the sample mean.

\end{tcolorbox}

\subsection{Large Sample Distribution of OLS
Estimator}\label{large-sample-distribution-of-ols-estimator}

Now we can establish the large sample distribution of the OLS estimator
\(\hat{\bbeta}\) based on CLT.

Under the Gauss-Markov assumptions, as \(n\to\infty\), \[
\sqrt{n} (\hat{\bbeta} - \bbeta) \xrightarrow{d} N(0, \sigma^2 E(\bX_i\bX_i')^{-1})
\]

\begin{itemize}
\item
  Again, note the scaling by \(\sqrt{n}.\) It shows how fast the
  estimator \(\hat{\bbeta}\) converges to the true parameter \(\bbeta\)
  as sample size \(n\) increases.
\item
  Note the difference of the variance between the exact distribution and
  the large sample distribution.

  \begin{itemize}
  \tightlist
  \item
    The exact distribution depends on the \textbf{sample data}
    \(\bX'\bX\)
  \item
    The large sample distribution depends on the \textbf{population
    moment} \(E(\bX_i\bX_i')\).
  \end{itemize}
\item
  It suggests that \(\hat{\bbeta}\) is \textbf{asymptotically normal}
  and \textbf{consistent}.

  ``Consistency'' means that as the sample size \(n\) increases, the
  estimator \(\hat{\bbeta}\) converges in probability to the true
  parameter \(\bbeta\).
\end{itemize}

\section{Measure of Fit}\label{measure-of-fit}

We often use \(R^2\) to measure the goodness of fit of a regression
model.

\begin{itemize}
\tightlist
\item
  \(R^2\) is called ``the coefficient of determination.''
\item
  It represents the proportion of variance in the dependent variable
  that is predictable by the regression model.
\item
  \(0\le R^2\le 1\)
\end{itemize}

\(R^2\) is defined as:

\[
R^2 = 1 - \frac{SSR}{SST} = \frac{SST - SSR}{SST} = \frac{SSE}{SST}
\]

where

\begin{itemize}
\tightlist
\item
  \(SSR = \sum_{i=1}^{n} (Y_i - \hat{Y}_i)^2 = \sum_{i=1}^{n} e_i^2\) is
  the sum of squared residuals (unexplained variation)
\item
  \(SST = \sum_{i=1}^{n} (Y_i - \bar{Y})^2\) is the total sum of squares
  (total variation)
\item
  \(SSE = \sum_{i=1}^{n} (\hat{Y}_i - \bar{Y})^2\) is the explained sum
  of squares (explained variation)
\item
  \(\bar{Y}= \frac{1}{n}\sum_{i=1}^n Y_i\) is the mean of the dependent
  variable
\end{itemize}

Note that adding variables always increases \(R^2\), even if the new
variables are not statistically significant. → To address this, we use
the \textbf{adjusted \(R^2\)}.

In a regression model with multiple explanatory variables, we often
use~\textbf{adjusted}~\(R^2\)~that adjusts the number of explanatory
variables.

``adjusted \(R^2\)'' or ``R-bar-squared'', defined by \[
\overline{R}^2 = 1- \frac{n-1}{n-k}(1-R^2) ,
\] which imposes a penalty as \(k\) increases in a given sample size
\(n\).

\section{Homoskedasticity}\label{homoskedasticity}

Homoskedasticity assumption requires constant variance of the error term
across all levels of the independent variable(s).

\[
\var(\varepsilon_i | X) = \sigma^2
\]

When the assumption is violated, we have heteroskedasticity:

\[
\var(\varepsilon_i | X) = \sigma^2_i
\]

that is, the variance of the error term varies with the level of the
independent variable(s).

If the error has heteroskedasticity, {the standard error under
homoskedasticity assumption will be underestimated}, leading to invalid
t statistics, affecting hypothesis tests and confidence intervals.

\subsection{Detection of
Heteroskedasticity}\label{detection-of-heteroskedasticity}

\subsection{Graphical Method}\label{graphical-method}

A common graphical method to detect heteroskedasticity is to plot the
residuals against the fitted values or one of the independent variables.

Recall the CASchools example, we regress TestScore on expenditure.

\begin{center}
\pandocbounded{\includegraphics[keepaspectratio]{04_Theoretical_Framework_LR_files/figure-pdf/unnamed-chunk-1-1.pdf}}
\end{center}

We recap the figure summary below:

\begin{itemize}
\item
  The residuals appear reasonably balanced around zero, though the
  spread suggests substantial unexplained variation.
\item
  The spread of residuals seems to \textbf{increase slightly with fitted
  values}, indicating potential heteroskedasticity (non-constant
  variance of errors). This suggests that the assumption of
  homoskedasticity may be violated, which could affect the reliability
  of our coefficient estimates standard errors.
\end{itemize}

\subsection{Fix Heteroskedasticity}\label{fix-heteroskedasticity}

A quick fix is to use {heteroskedasticity-robust standard errors}, which
are valid even when the homoskedasticity assumption is violated.

\begin{Shaded}
\begin{Highlighting}[]
\FunctionTok{library}\NormalTok{(tidyverse)}
\FunctionTok{library}\NormalTok{(sandwich) }\CommentTok{\# for robust standard errors}
\FunctionTok{library}\NormalTok{(lmtest) }\CommentTok{\# for coeftest()}
\FunctionTok{library}\NormalTok{(stargazer)}

\NormalTok{f\_name }\OtherTok{\textless{}{-}} \StringTok{"https://raw.githubusercontent.com/my1396/FIN5005{-}Fall2025/refs/heads/main/data/CASchools\_test\_score.csv"}
\NormalTok{cas }\OtherTok{\textless{}{-}} \FunctionTok{read\_csv}\NormalTok{(f\_name,}
    \AttributeTok{col\_types =} \FunctionTok{cols}\NormalTok{(}
        \AttributeTok{county =} \FunctionTok{col\_factor}\NormalTok{(), }\CommentTok{\# read as factor}
        \AttributeTok{grades =} \FunctionTok{col\_factor}\NormalTok{()}
\NormalTok{    )}
\NormalTok{)}
\NormalTok{cas }\OtherTok{\textless{}{-}}\NormalTok{ cas }\SpecialCharTok{\%\textgreater{}\%}
    \FunctionTok{mutate}\NormalTok{(}\AttributeTok{TestScore =}\NormalTok{ (read }\SpecialCharTok{+}\NormalTok{ math) }\SpecialCharTok{/} \DecValTok{2}\NormalTok{)}
\NormalTok{model }\OtherTok{\textless{}{-}} \FunctionTok{lm}\NormalTok{(TestScore }\SpecialCharTok{\textasciitilde{}}\NormalTok{ expenditure, }\AttributeTok{data =}\NormalTok{ cas)}
\end{Highlighting}
\end{Shaded}

\begin{Shaded}
\begin{Highlighting}[]
\CommentTok{\# OLS standard errors}
\NormalTok{se\_ols }\OtherTok{\textless{}{-}} \FunctionTok{sqrt}\NormalTok{(}\FunctionTok{diag}\NormalTok{(}\FunctionTok{vcov}\NormalTok{(model)))}

\CommentTok{\# HC1 heteroskedasticity{-}robust standard errors}
\NormalTok{se\_hc1 }\OtherTok{\textless{}{-}} \FunctionTok{sqrt}\NormalTok{(}\FunctionTok{diag}\NormalTok{(}\FunctionTok{vcovHC}\NormalTok{(model, }\AttributeTok{type =} \StringTok{"HC1"}\NormalTok{)))}

\CommentTok{\# Stargazer output: same model, different SEs}
\FunctionTok{stargazer}\NormalTok{(model, model,}
          \AttributeTok{se =} \FunctionTok{list}\NormalTok{(se\_ols, se\_hc1),}
          \AttributeTok{digits =} \DecValTok{4}\NormalTok{,}
          \AttributeTok{column.labels =} \FunctionTok{c}\NormalTok{(}\StringTok{"OLS SE"}\NormalTok{, }\StringTok{"Robust SE"}\NormalTok{),}
          \AttributeTok{type =} \StringTok{"text"}
\NormalTok{          )}
\end{Highlighting}
\end{Shaded}

\begin{verbatim}

===========================================================
                                   Dependent variable:     
                               ----------------------------
                                        TestScore          
                                   OLS SE       Robust SE  
                                    (1)            (2)     
-----------------------------------------------------------
expenditure                      0.0057***      0.0057***  
                                  (0.0014)      (0.0016)   
                                                           
Constant                        623.6165***    623.6165*** 
                                  (7.7197)      (8.4664)   
                                                           
-----------------------------------------------------------
Observations                        420            420     
R2                                 0.0366        0.0366    
Adjusted R2                        0.0343        0.0343    
Residual Std. Error (df = 418)    18.7239        18.7239   
F Statistic (df = 1; 418)        15.8734***    15.8734***  
===========================================================
Note:                           *p<0.1; **p<0.05; ***p<0.01
\end{verbatim}

\textbf{Under the heteroskedasticity-robust standard errors:}

\begin{itemize}
\tightlist
\item
  No effects on the coefficient estimates.
\item
  The Std. Error is slightly larger (0.001620 vs.~0.001443). The
  coefficient remains statistically significant at the 5\% level.
\end{itemize}

\textbf{Interpretation}

The heteroskedasticity-robust results suggest that the OLS standard
errors may have been {\textbf{too optimistic (underestimated
variability)}}. While the statistical significance of the expenditure
effect is slightly reduced, it remains highly significant, indicating a
robust positive relationship between expenditure and test scores.


\backmatter


\end{document}
